% !TEX root = ../metatime_monograph.tex
\chapter{Higgs Vacuum Expectation Value}
\label{ch:higgs_vev}


% CITATION_CONTEXT_V10_9_BEGIN
\paragraph{Established context.}
Spontaneous electroweak symmetry breaking and gauge-boson mass generation are
established by the Brout--Englert, Higgs, and Guralnik--Hagen--Kibble mechanisms
\cite{EnglertBrout1964,Higgs1964,GHK1964}.  Electroweak normalization follows
Glashow--Weinberg--Salam and modern reference data
\cite{Glashow1961,Weinberg1967,Salam1968,PDG2024,CODATA2022}; the Metatime
expression for the vacuum expectation value is an additional derivation claim.
% CITATION_CONTEXT_V10_9_END

\section{Derivation}

The electroweak scale is encoded in baryonic quantum numbers.  The Higgs
vacuum expectation value is:

\begin{equation}
v = E_{\text{proton}}\,(L_3^2 L_5 + L_3 L_4 + L_4).
\label{eq:vev}
\end{equation}

\subsection{Term-by-Term Breakdown}

\[
\begin{aligned}
L_3^2 L_5 &= 7^2 \times 5 = 49 \times 5 = 245, \\
L_3 L_4   &= 7 \times 2 = 14, \\
L_4       &= 2.
\end{aligned}
\]

Sum: \(245 + 14 + 2 = 261\).

\subsection{Numerical Evaluation}

\[
v = 938.272 \times 261\ \mev = 244.89\ \gev.
\]

\section{Comparison with PDG}

The PDG value is \(v^{\text{PDG}} = 246.22\)~GeV.

\[
\frac{244.89 - 246.22}{246.22} = -0.54\%.
\]

\section{Physical Interpretation}

The VEV expression has a clear structural interpretation:
\begin{itemize}
\item \(L_3^2 L_5 = 245\) represents the SU(3) colour factor squared times
the intention dimension.
\item \(L_3 L_4 = 14\) is the colour--annihilation mixing.
\item \(L_4 = 2\) is the annihilation constant itself, representing the
boundary ``thickness'' of the Poincar\'{e} disk.
\end{itemize}
